% GENERATED FILE. Edit canonical lessons, then run npm run generate:books.
% canonical-source-hash: fnv1a64:080c4a6b9e210c27
% canonical-lessons: IT-C450-sagra, IT-C450-corteo, IT-C450-processione, IT-C450-patrono, IT-C450-bancarella

\chapter{Feste di paese}
\label{ch:it-feste-di-paese}
\hlchaptermodality{450}

\begin{chapteropening}
\textbf{By the end of this chapter:} \emph{I can talk about a village festival: the food festival, the parade, the procession, the patron saint and the stalls.}

Complete the last lesson of chapter 450: I can talk about a village festival: the food festival, the parade, the procession, the patron saint and the stalls.
\end{chapteropening}

\section[la sagra]{la sagra — the village food festival}

\label{lesson:IT-C450-sagra}



\begin{quote}
Before the new one: say the Italian for to collapse, then the Italian for to clear.
\end{quote}

\subsection*{You'll want to know: la sagra}
\textbf{la sagra} — \textquotedblleft{}the village food festival\textquotedblright{}.

\begin{quote}
Domenica c'è la sagra del pesce.
\end{quote}

\begin{quote}
\emph{On Sunday there is the fish festival.}
\end{quote}

Each one celebrates a local food: chestnuts, mushrooms, wine.

\begin{cousinweb}
\emph{sagra} is inherited from Latin \emph{sacrum} 'sacred rite' (EtymDB), from \emph{sacer} 'holy': first the feast of a church.
\end{cousinweb}

\begin{grammarlens}[title={What you've built}]
One more word for telling, explaining and describing at B1.
\end{grammarlens}

\subsection*{Your turn}
\begin{itemize}
\raggedright
  \item \emph{Say it:} \emph{la sagra}
  \item \emph{Say it:} \emph{la sagra}, once more
  \item \emph{Say it:} a sentence of your own with \emph{la sagra}
  \item \emph{Recall:} say the Italian for to collapse, then the Italian for to clear, then say \emph{la sagra} again
\end{itemize}

\begin{tcolorbox}[breakable,colback=teal!4,colframe=teal!35!black,title={Before you move on}]
What does la sagra mean? (\textquotedblleft{}The village food festival\textquotedblright{}.) Say it once more.
\end{tcolorbox}
\section[il corteo]{il corteo — the parade, the procession, the march}

\label{lesson:IT-C450-corteo}



\begin{quote}
Before the new one: say the Italian for to clear, then the Italian for the village food festival.
\end{quote}

\subsection*{You'll want to know: il corteo}
\textbf{il corteo} — \textquotedblleft{}the parade, the procession, the march\textquotedblright{}.

\begin{quote}
Il corteo passa per la via principale.
\end{quote}

\begin{quote}
\emph{The parade goes along the main street.}
\end{quote}

A political march too: \emph{un corteo di protesta}.

\begin{cousinweb}
\emph{corteo} is built on \emph{corte} 'court', which is inherited from Latin \emph{cohors, cohortis} 'yard, retinue' (Lewis \& Short): first the people following a lord.
\end{cousinweb}

\begin{grammarlens}[title={What you've built}]
One more word for telling, explaining and describing at B1.
\end{grammarlens}

\subsection*{Your turn}
\begin{itemize}
\raggedright
  \item \emph{Say it:} \emph{il corteo}
  \item \emph{Say it:} \emph{il corteo}, once more
  \item \emph{Say it:} a sentence of your own with \emph{il corteo}
  \item \emph{Recall:} say the Italian for to clear, then the Italian for the village food festival, then say \emph{il corteo} again
\end{itemize}

\begin{tcolorbox}[breakable,colback=teal!4,colframe=teal!35!black,title={Before you move on}]
What does il corteo mean? (\textquotedblleft{}The parade, the procession, the march\textquotedblright{}.) Say it once more.
\end{tcolorbox}
\section[la processione]{la processione — the procession (religious)}

\label{lesson:IT-C450-processione}



\begin{quote}
Before the new one: say the Italian for the village food festival, then the Italian for the parade.
\end{quote}

\subsection*{You'll want to know: la processione}
\textbf{la processione} — \textquotedblleft{}the procession (religious)\textquotedblright{}.

\begin{quote}
La processione esce dalla chiesa alle sei.
\end{quote}

\begin{quote}
\emph{The procession leaves the church at six.}
\end{quote}

It often follows a statue of the patron saint.

\begin{cousinweb}
\emph{processione} is a learned borrowing from Latin \emph{prōcessiō} 'a going forward' (EtymDB), from \emph{prōcēdere} 'to go forward', built on \emph{cēdere} 'to go'.
\end{cousinweb}

\begin{grammarlens}[title={What you've built}]
One more word for telling, explaining and describing at B1.
\end{grammarlens}

\subsection*{Your turn}
\begin{itemize}
\raggedright
  \item \emph{Say it:} \emph{la processione}
  \item \emph{Say it:} \emph{la processione}, once more
  \item \emph{Say it:} a sentence of your own with \emph{la processione}
  \item \emph{Recall:} say the Italian for the village food festival, then the Italian for the parade, then say \emph{la processione} again
\end{itemize}

\begin{tcolorbox}[breakable,colback=teal!4,colframe=teal!35!black,title={Before you move on}]
What does la processione mean? (\textquotedblleft{}The procession (religious)\textquotedblright{}.) Say it once more.
\end{tcolorbox}
\section[il patrono]{il patrono — the patron saint; the patron}

\label{lesson:IT-C450-patrono}



\begin{quote}
Before the new one: say the Italian for the parade, then the Italian for the procession.
\end{quote}

\subsection*{You'll want to know: il patrono}
\textbf{il patrono} — \textquotedblleft{}the patron saint; the patron\textquotedblright{}.

\begin{quote}
Il giorno del patrono i negozi sono chiusi.
\end{quote}

\begin{quote}
\emph{On the patron saint's day the shops are closed.}
\end{quote}

\emph{La festa patronale} — the patron saint's feast.

\begin{cousinweb}
\emph{patrono} is a learned borrowing from Latin \emph{patrōnus} 'protector, patron' (EtymDB), from \emph{pater} 'father'.
\end{cousinweb}

\begin{grammarlens}[title={What you've built}]
One more word for telling, explaining and describing at B1.
\end{grammarlens}

\subsection*{Your turn}
\begin{itemize}
\raggedright
  \item \emph{Say it:} \emph{il patrono}
  \item \emph{Say it:} \emph{il patrono}, once more
  \item \emph{Say it:} a sentence of your own with \emph{il patrono}
  \item \emph{Recall:} say the Italian for the parade, then the Italian for the procession, then say \emph{il patrono} again
\end{itemize}

\begin{tcolorbox}[breakable,colback=teal!4,colframe=teal!35!black,title={Before you move on}]
What does il patrono mean? (\textquotedblleft{}The patron saint; the patron\textquotedblright{}.) Say it once more.
\end{tcolorbox}
\section[la bancarella]{la bancarella — the stall (at a market)}

\label{lesson:IT-C450-bancarella}



\begin{quote}
Before the new one: say the Italian for the procession, then the Italian for the patron saint.
\end{quote}

\subsection*{You'll want to know: la bancarella}
\textbf{la bancarella} — \textquotedblleft{}the stall (at a market)\textquotedblright{}.

\begin{quote}
Alle bancarelle vendono dolci e giocattoli.
\end{quote}

\begin{quote}
\emph{At the stalls they sell sweets and toys.}
\end{quote}

Built on \emph{banco}: a little bench to sell from.

\begin{cousinweb}
\emph{bancarella} is built on \emph{banco} 'bench, counter', which was borrowed from a Germanic word for a bench (EtymDB).
\end{cousinweb}

\begin{grammarlens}[title={What you've built}]
One more word for telling, explaining and describing at B1.
\end{grammarlens}

\subsection*{Your turn}
\begin{itemize}
\raggedright
  \item \emph{Say it:} \emph{la bancarella}
  \item \emph{Say it:} \emph{la bancarella}, once more
  \item \emph{Say it:} a sentence of your own with \emph{la bancarella}
  \item \emph{Recall:} say the Italian for the procession, then the Italian for the patron saint, then say \emph{la bancarella} again
\end{itemize}

\begin{tcolorbox}[breakable,colback=teal!4,colframe=teal!35!black,title={Before you move on}]
What does la bancarella mean? (\textquotedblleft{}The stall (at a market)\textquotedblright{}.) Say it once more.
\end{tcolorbox}
